\chapter{Data Pipeline}
\label{ch:data}

BabyGPT trains on a single 1\,MB text file that fits in RAM.  GPT-2 was trained on WebText, a 40\,GB corpus of web pages.  Modern reproductions use even larger, openly available datasets like FineWeb-Edu.  In this chapter we build a data pipeline that downloads shards in Parquet format, tokenizes documents with BOS prepended, packs them efficiently into fixed-length sequences using best-fit bin packing, and distributes batches across multiple GPUs.


\section{Web-Scale Datasets}

The original GPT-2 was trained on \textbf{WebText}, a custom dataset of web pages filtered by Reddit upvotes (links with $\geq 3$ karma).  WebText was never publicly released.  Several open alternatives now exist:

\begin{center}
\begin{tabular}{lrl}
\toprule
Dataset & Size & Description \\
\midrule
The Pile & 800 GB & 22 diverse sources (academic, books, code, web) \\
C4 & 750 GB & Cleaned Common Crawl \\
FineWeb & 15 TB & Cleaned Common Crawl (2013--2024) \\
FineWeb-Edu & 1.3 TB & FineWeb filtered for educational content \\
ClimbMix & 400B tokens & Nanochat's curated multi-source blend \\
\bottomrule
\end{tabular}
\end{center}

We use \textbf{FineWeb-Edu}: it is openly available on HuggingFace, well-documented, and filtered for quality.  For a GPT-2 scale reproduction ($\sim$10B tokens), we download a small subset (10--50 shards, each $\sim$100M tokens).


\section{The Parquet Format}

FineWeb-Edu is distributed as Apache Parquet files.  Parquet is a columnar storage format designed for large-scale data:

\begin{itemize}
  \item \textbf{Columnar}: data is stored column-by-column, so reading just the \code{text} column skips all metadata columns.
  \item \textbf{Compressed}: Snappy or Zstandard compression reduces file sizes by 2--5x.
  \item \textbf{Row groups}: each file is divided into row groups (typically 10K--100K rows each).  You can read one row group at a time without loading the entire file.
\end{itemize}

The \code{pyarrow} library provides fast Parquet reading in Python.


\section{Dataset Management}

The \code{microgpt/dataset.py} module handles downloading and listing parquet shards:

\filetag{microgpt/dataset.py}
\begin{lstlisting}
"""Dataset management: download and list FineWeb-Edu parquet shards.

FineWeb-Edu is an open, high-quality web text dataset filtered for
educational content.  It is distributed as parquet files on HuggingFace.
Each file contains a 'text' column with document strings.

Usage:
    python -m microgpt.dataset                  # download 10 shards (~5 GB)
    python -m microgpt.dataset --num-shards 50  # download more for full training
"""

import os
import glob

import requests


# ---------------------------------------------------------------------------
# Configuration
# ---------------------------------------------------------------------------

DATA_DIR = os.path.join(os.path.dirname(os.path.dirname(__file__)), "data", "microgpt")

# FineWeb-Edu-score-2 sample: ~10B tokens total, split into 100 parquet files
# Each shard is ~50-100 MB, containing ~100M tokens
BASE_URL = (
    "https://huggingface.co/datasets/HuggingFaceFW/fineweb-edu-score-2"
    "/resolve/main/data/CC-MAIN-2024-10"
)
# Shard filenames: train-00000-of-00099.parquet ... train-00098-of-00099.parquet
NUM_SHARDS_AVAILABLE = 99


def shard_filename(idx: int) -> str:
    return f"train-{idx:05d}-of-{NUM_SHARDS_AVAILABLE:05d}.parquet"


def shard_path(idx: int) -> str:
    return os.path.join(DATA_DIR, shard_filename(idx))


# ---------------------------------------------------------------------------
# Download
# ---------------------------------------------------------------------------

def download_shard(idx: int) -> None:
    """Download a single parquet shard if not already present."""
    path = shard_path(idx)
    if os.path.exists(path):
        return
    url = f"{BASE_URL}/{shard_filename(idx)}"
    print(f"  Downloading shard {idx}: {url}")
    os.makedirs(DATA_DIR, exist_ok=True)
    response = requests.get(url, stream=True, timeout=60)
    response.raise_for_status()
    # Write to a temp file first, then rename (atomic on POSIX)
    tmp_path = path + ".tmp"
    with open(tmp_path, "wb") as f:
        for chunk in response.iter_content(chunk_size=8 * 1024 * 1024):
            f.write(chunk)
    os.rename(tmp_path, path)


def download_shards(num_shards: int = 10) -> None:
    """Download parquet shards for training."""
    num_shards = min(num_shards, NUM_SHARDS_AVAILABLE)
    print(f"Downloading {num_shards} FineWeb-Edu shards to {DATA_DIR}")
    for idx in range(num_shards):
        download_shard(idx)
    print(f"Done. {num_shards} shards in {DATA_DIR}")


# ---------------------------------------------------------------------------
# Listing
# ---------------------------------------------------------------------------

def list_parquet_files() -> list[str]:
    """Return sorted list of all downloaded parquet shard paths."""
    pattern = os.path.join(DATA_DIR, "train-*.parquet")
    paths = sorted(glob.glob(pattern))
    return paths
\end{lstlisting}

\checkpoint{Download two shards and verify:}

\shelltag
\begin{lstlisting}[language=bash]
uv run python -m microgpt.dataset --num-shards 2
\end{lstlisting}


\section{BOS-Aligned Best-Fit Packing}

\subsection{The packing problem}

Documents have variable lengths.  Training requires fixed-length sequences of $T+1$ tokens (input + shifted target).  How do we fill each row?

\textbf{Naive concatenation} glues documents end-to-end and chops the result into $T$-length chunks.  This is simple and wastes nothing, but a model token at position 500 might attend back to a BOS token at position 200 that belongs to a completely different document.  The attention is wasted on irrelevant context.

\textbf{BOS-aligned packing} ensures every row starts with a BOS token and that each token can attend back to its own document's beginning.  The tradeoff: documents that don't fit a row are cropped, wasting $\sim$35\% of tokens.

\subsection{The best-fit algorithm}

For each row in the batch:
\begin{enumerate}
  \item Maintain a buffer of tokenized documents.
  \item Find the \textbf{largest} document that fits entirely in the remaining space.  Append it to the row.  Repeat.
  \item When no document fits, \textbf{crop} the shortest document in the buffer to fill the remaining space exactly.
\end{enumerate}

This achieves 100\% utilization (no padding) while minimizing the amount of cropped text by preferring documents that fit without cropping.


\section{The Distributed Dataloader}

\filetag{microgpt/dataloader.py}
\begin{lstlisting}
"""Distributed dataloader with BOS-aligned best-fit packing.

This dataloader:
1. Reads documents from parquet shards (via dataset.py)
2. Tokenizes them with BOS prepended
3. Packs documents into fixed-length rows using best-fit bin packing
4. Distributes shards across DDP ranks
5. Pre-allocates pinned CPU buffers for efficient GPU transfer

Every row starts with a BOS token, so the model always sees the beginning
of at least one document.  Documents that don't fit are cropped (the cropped
portion is discarded, not carried over).  This wastes ~35% of tokens but
ensures every token can attend back to its document's BOS.
"""

import torch
import pyarrow.parquet as pq

from .dataset import list_parquet_files


# ---------------------------------------------------------------------------
# Document iterator (with DDP sharding)
# ---------------------------------------------------------------------------

def _document_batches(
    split: str,
    rank: int = 0,
    world_size: int = 1,
    batch_size: int = 128,
):
    """Infinite iterator over batches of document strings from parquet files.

    DDP sharding: each rank reads every world_size-th row group, so ranks
    see disjoint data without any communication.

    Yields: (text_batch, shard_idx)
    """
    paths = list_parquet_files()
    assert paths, "No parquet files found. Run: python -m microgpt.dataset"
    # Last shard is validation, rest are training
    if split == "train":
        paths = paths[:-1]
    else:
        paths = paths[-1:]

    while True:  # infinite iteration (multi-epoch)
        for shard_idx, filepath in enumerate(paths):
            pf = pq.ParquetFile(filepath)
            # Each rank reads every world_size-th row group
            rg_idx = rank
            while rg_idx < pf.num_row_groups:
                rg = pf.read_row_group(rg_idx)
                texts = rg.column("text").to_pylist()
                for i in range(0, len(texts), batch_size):
                    yield texts[i : i + batch_size], shard_idx
                rg_idx += world_size
\end{lstlisting}

DDP sharding is simple: each rank reads every $N$\textsuperscript{th} row group (where $N$ is the world size).  Rank 0 reads row groups 0, $N$, $2N$, \ldots; rank 1 reads 1, $N+1$, $2N+1$, \ldots.  No communication is needed --- each rank sees disjoint data.

\filetag{microgpt/dataloader.py}
\begin{lstlisting}
def _pack_row(
    doc_buffer: list[list[int]],
    row: torch.Tensor,
    row_capacity: int,
) -> None:
    """Pack documents from doc_buffer into a single row using best-fit.

    Algorithm:
    1. Find the largest document that fits entirely in the remaining space.
    2. Append it and repeat.
    3. When no document fits, crop the shortest one to fill exactly.

    This achieves 100% utilization (no padding) at the cost of ~35% cropped
    tokens.
    """
    pos = 0
    while pos < row_capacity:
        remaining = row_capacity - pos

        # Find largest doc that fits entirely
        best_idx = -1
        best_len = 0
        for i, doc in enumerate(doc_buffer):
            doc_len = len(doc)
            if doc_len <= remaining and doc_len > best_len:
                best_idx = i
                best_len = doc_len

        if best_idx >= 0:
            doc = doc_buffer.pop(best_idx)
            row[pos : pos + len(doc)] = torch.tensor(doc, dtype=torch.long)
            pos += len(doc)
        else:
            # No doc fits --- crop shortest to fill remaining space exactly
            shortest_idx = min(
                range(len(doc_buffer)), key=lambda i: len(doc_buffer[i])
            )
            doc = doc_buffer.pop(shortest_idx)
            row[pos : pos + remaining] = torch.tensor(
                doc[:remaining], dtype=torch.long
            )
            pos += remaining
\end{lstlisting}

\filetag{microgpt/dataloader.py}
\begin{lstlisting}
def create_dataloader(
    tokenizer,
    batch_size: int,
    sequence_len: int,
    split: str = "train",
    rank: int = 0,
    world_size: int = 1,
    device: str = "cuda",
    buffer_size: int = 1000,
):
    """Create an infinite dataloader that yields (inputs, targets) batches.

    Args:
        tokenizer: BPETokenizer instance (needs .encode_with_bos and .bos_id)
        batch_size: number of rows per batch (per GPU)
        sequence_len: context length (T)
        split: "train" or "val"
        rank: DDP rank (0 for single-GPU)
        world_size: DDP world size (1 for single-GPU)
        device: target device ("cuda" or "cpu")
        buffer_size: min documents in buffer before packing

    Yields:
        (inputs, targets): both shape (batch_size, sequence_len), on device
    """
    row_capacity = sequence_len + 1  # +1 because targets are shifted by 1
    batches = _document_batches(split, rank, world_size)
    doc_buffer: list[list[int]] = []

    # Pre-allocate buffers for efficient transfer
    use_cuda = device == "cuda"
    row_buffer = torch.empty(
        (batch_size, row_capacity), dtype=torch.long
    )
    cpu_buffer = torch.empty(
        2 * batch_size * sequence_len, dtype=torch.long, pin_memory=use_cuda
    )
    gpu_buffer = torch.empty(
        2 * batch_size * sequence_len, dtype=torch.long, device=device
    )
    # Views into the flat buffers
    cpu_inputs = cpu_buffer[: batch_size * sequence_len].view(
        batch_size, sequence_len
    )
    cpu_targets = cpu_buffer[batch_size * sequence_len :].view(
        batch_size, sequence_len
    )
    inputs = gpu_buffer[: batch_size * sequence_len].view(
        batch_size, sequence_len
    )
    targets = gpu_buffer[batch_size * sequence_len :].view(
        batch_size, sequence_len
    )

    while True:
        # Build a full batch
        for row_idx in range(batch_size):
            # Refill buffer if needed
            while len(doc_buffer) < buffer_size:
                text_batch, _ = next(batches)
                for text in text_batch:
                    doc_ids = tokenizer.encode_with_bos(text)
                    doc_buffer.append(doc_ids)

            _pack_row(doc_buffer, row_buffer[row_idx], row_capacity)

        # Split into inputs and targets (shifted by 1)
        cpu_inputs.copy_(row_buffer[:, :-1])
        cpu_targets.copy_(row_buffer[:, 1:])

        # Single host-to-device transfer
        gpu_buffer.copy_(cpu_buffer, non_blocking=use_cuda)
        yield inputs, targets
\end{lstlisting}

\subsection{Buffer management}

The dataloader maintains a \textbf{document buffer} of at least \code{buffer\_size} tokenized documents.  When the buffer runs low, it reads the next batch of raw text from parquet, tokenizes it (prepending BOS), and appends the token lists to the buffer.

Pre-allocated pinned CPU buffers and a single host-to-device copy per batch minimize transfer overhead.

\tip{Why is the row capacity $T+1$ instead of $T$?  Because the training targets are the inputs shifted by one position.  A row of $T+1$ tokens gives $T$ (input, target) pairs: input = row[0:T], target = row[1:T+1].  This is the same teacher-forcing setup as BabyGPT, just with more careful buffer management.}


\section{Summary}

The data pipeline has three components:

\begin{enumerate}
  \item \textbf{Dataset manager} (\code{microgpt/dataset.py}): downloads parquet shards from HuggingFace and lists them.
  \item \textbf{Document iterator}: reads row groups from parquet files, sharded across DDP ranks, with infinite looping.
  \item \textbf{BOS-aligned dataloader} (\code{microgpt/dataloader.py}): tokenizes documents, packs them into fixed-length rows using best-fit bin packing, and transfers batches to GPU.
\end{enumerate}

With the data pipeline in place, the next chapter sets up the distributed training infrastructure (DDP, mixed precision, gradient accumulation) needed to actually train on this data.
